\newcommand\llbBra{\bullet\{\!\!\{ }
\newcommand\llBra{\{\!\!\{}
\newcommand\rrBra{\}\!\!\}}

%\newcommand{\reach}[1]{\mathcal{R}(#1)}
\newcommand{\reach}[1]{{\it reach}(#1)}
\newcommand{\ceq}{\asymp}
\newcommand{\rt}{\leftharpoondown\rightharpoonup}
\newcommand{\source}{\mathtt{source}}
\newcommand{\target}{\mathtt{target}}
\newcommand{\ft}{\mathtt{t}}
\newcommand{\reverse}[1]{\leftharpoondown{#1}}
\newcommand{\len}[1]{|#1|}
\newcommand{\fw}[1]{\triangleright{#1}}
\newcommand{\bw}[1]{\triangleleft{#1}}
\newcommand{\afw}{\triangleright\!\triangleright}
\newcommand{\abw}{\triangleleft\!\triangleleft}

\newcommand{\m}{Mazurkiewicz }

%
% FIGURE SEPARATOR
%
%\newcommand{\topfigrule}{\vskip3pt\noindent\rule{\textwidth}{1pt}\vskip-15pt}

%
% MACRO FOR COMMENTS
%
\newcommand{\nota}[1]{\noindent \fbox{ \parbox{\textwidth}{#1} }  }

%
% ABBREVIATIONS AND SYMBOLS
%

%\newcommand{\eg}{e.g.}
\newcommand{\ie}{i.e.}
\newcommand{\wrt}{w.r.t.}

%
% MULTISETS
%
\newcommand{\supp}{\mathit{supp}}

%
% NETS
%
\newcommand{\minp}[1]{\ensuremath{{}^\circ{#1}}}
\newcommand{\maxp}[1]{\ensuremath{#1^\circ}}

\newcommand{\preS}[1]{\ensuremath{{}^\bullet{#1}}}
\newcommand{\postS}[1]{\ensuremath{#1^\bullet}}

\newcommand{\unf}{\mathcal{U}}
\newcommand{\places}{\alpha}
\newcommand{\p}[1]{{\sf #1}}

\newcommand{\transitions}{\beta}
\newcommand{\tr}[1]{{\sf #1}}

\newcommand{\nat}{\mathbb{N}}
%\newcommand{\nat}{\bbbn}

\newcommand{\en}[1]{\textsf{en}(#1)}
\newcommand{\co}[2]{#1 \; co \; #2}

\newcommand{\multiset}[1]{#1}

%
% REDUCTIONS
%
%\newcommand{\red}{\rightarrow}
% \newcommandx{\red}[1]{\xrightarrow{#1}}
%\newcommand{\fred}{\mapsto}
% \makeatletter
% \newcommand{\xMapsto}[2][]{\ext@arrow 0599{\Mapstofill@}{#1}{#2}}
% \def\Mapstofill@{\arrowfill@{\Mapstochar\Relbar}\Relbar\Rightarrow}
% \makeatother
%
%\newcommandx{\fred}[1][1=\tr t]{\xmapsto{#1}}

% \newcommandx{\fred}[1]{\stackrel{#1}{\twoheadrightarrow}}
%
% \newcommandx{\bred}[1]{\stackrel{#1}{\rightsquigarrow}}
% \newcommand{\rred}{\Rightarrow}
% \newcommand{\jeq}{\; \mathop{\triangleright} \;}
%
\newcommand{\pntrans}{\llbracket\rangle}
\newcommand{\exec}[1]{\llbracket#1\rangle}
%
% INFERENCE RULES
%
% \def \mathrule #1#2#3{\begin{array}{l}%
%     {\mbox{\scriptsize ({\sc #1})} }%
%     \\ \irule{#2}{#3}%
% \end{array}}
% \newcommand{\irule}[2]{\frac{\textstyle\rule[-1.3ex]{0cm}{3ex}#1}%
% {\textstyle\rule[-.5ex]{0cm}{3ex}#2}}
% \newcommand{\arule}[3]{\frac{\textstyle\rule[-.8ex]{0cm}{3ex}#1}%
% {\textstyle\rule[.2ex]{0cm}{3ex}#2}{\mbox{\scriptsize {\sc #3}}}}
%
% %
% % COLOUR-NETS
% %
%
% \newcommand{\subsc}{\star}
%
% \makeatletter
% \newsavebox{\@brx}
% \newcommand{\llangle}[1][]{\savebox{\@brx}{\(\m@th{#1\langle}\)}%
%   \mathopen{\copy\@brx\mkern2mu\kern-0.9\wd\@brx\usebox{\@brx}}}
% \newcommand{\rrangle}[1][]{\savebox{\@brx}{\(\m@th{#1\rangle}\)}%
%   \mathclose{\copy\@brx\mkern2mu\kern-0.9\wd\@brx\usebox{\@brx}}}
% \makeatother
%
\newcommand{\enabled}[1]{\llbracket #1\rrbracket}
\newcommand{\fir}[1]{\llbracket #1 \rangle}
\newcommand{\fs}[1]{\llbracket #1 \rrangle}
\newcommand{\TT}{\mathsf{TransitionType}}
\newcommand{\fwd}{\mathsf{fwd}}
\newcommand{\bwd}{\mathsf{bwd}}
\newcommand{\Pl}{\mathsf{Pl}}
\newcommand{\Tr}{\mathsf{Tr}}
%
% REVERSIBLE
%

% \makeatletter
% \DeclareRobustCommand{\cev}[1]{%
%   \mathpalette\do@cev{#1}%
% }
% \newcommand{\do@cev}[2]{%
%   \fix@cev{#1}{+}%
%   \reflectbox{$\m@th#1\overrightarrow{\reflectbox{$\fix@cev{#1}{-}\m@th#1#2\fix@cev{#1}{+}$}}$}%
%   \fix@cev{#1}{-}%
% }
% \newcommand{\fix@cev}[2]{%
%   \ifx#1\displaystyle
%     \mkern#23mu
%   \else
%     \ifx#1\textstyle
%       \mkern#23mu
%     \else
%       \ifx#1\scriptstyle
%         \mkern#22mu
%       \else
%         \mkern#22mu
%       \fi
%     \fi
%   \fi
% }
%
% \makeatother
%
% \newcommand{\leadstoop}[1]{\stackrel{\reflectbox{\footnotesize$\leadsto$}}{#1}}
%
% %\newcommand{\rev}[1]{\cev{#1}}
%
\newcommand{\rev}[1]{\leadstoop{#1}}


%%%%Drawing %%%%

% \definecolor{mygray}{RGB}{120,120,120}

% \newcommand{\drawplace}{*[o]=<1.2pc,1.2pc>{\ }\drop\cir{}}
% \newcommand{\drawmarkedplace}{*[o]=<1.2pc,1.2pc>{\bullet}\drop\cir{}}
% \newcommand{\drawcolouredmarkedplace}[1]{*[o]=<1.2pc,1.2pc>{#1}\drop\cir{}}
%
% \newcommand{\nameplaceup}[1]{\POS[]+<0pc,1pc>\drop{{#1}}}
% \newcommand{\nameplacedown}[1]{\POS[]-<0pc,1pc>\drop{{#1}}}
% \newcommand{\nameplaceright}[1]{\POS[]+<1pc,0pc>\drop{{#1}}}
% \newcommand{\nameplaceleft}[1]{\POS[]-<1pc,0pc>\drop{{#1}}}
%
%
% \newcommand{\drawtrans}[1]{*=<2pc,1.3pc>{{#1}}\drop\frm{-}}
% \newcommand{\drawtransr}[1]{*=<2pc,1pc>[F:mygray]{#1}}
%
%
% %%% 1-inf-safe
% \newcommand{\OneInfSafe}{$1$-$\infty$-safe\xspace}
% \newcommand{\OneInfMarking}{$1$-$\infty$ marking\xspace}
%
% %%% coloured nets
% \newcommand{\colourset}{\mathcal{C}}
% \newcommand{\historyset}{\mathcal{H}}
% \newcommand{\varset}{\mathcal{X}}
% \newcommand{\vars}[1]{{\it vars}(#1)}
%

%--------------------------------------------------------------
%General Environments
%--------------------------------------------------------------
%--------------------------------------------------------------

\newcommand{\be}{\begin{enumerate}}
\newcommand{\ee}{\end{enumerate}}
\newcommand{\bd}{\begin{description}}
\newcommand{\ed}{\end{description}}
\newcommand{\bi}{\begin{itemize}}
\newcommand{\ei}{\end{itemize}}


%\newenvironment{proof}{\par \smallskip{\bf Proof:}}{\hfill\stopproof}
%\def\stopproof{\square}
% \def\square{\vbox{\hrule height.2pt\hbox{\vrule width.2pt height5pt \kern5pt
% \vrule width.2pt} \hrule height.2pt}}
%
%
% \newcommand{\tw}{\mathrm{tw}}
%
% %--------------------------------------------------------------
% %--------------------------------------------------------------
% %Figures and such
% %--------------------------------------------------------------
% %--------------------------------------------------------------
% \newcommand{\scale}[2]{\scalebox{#1}{#2}}
%
% \newcommand{\fig}[1]{
% \begin{figure}[h]
% \rotatebox{0}{\includegraphics{#1}}
% \end{figure}}
%
% \newcommand{\figcap}[2]
% {
% \begin{figure}[h]
% \rotatebox{270}{\includegraphics{#1}} \caption{#2}
% \end{figure}
% }
%
% \newcommand{\scalefig}[2]{
% \begin{figure}[h]
% \scalebox{#1}{\includegraphics{#2}}
% \end{figure}}
%
% \newcommand{\scalefigcap}[3]{
% \begin{figure}[h]
% \scalebox{#1}{\rotatebox{0}{\includegraphics{#2}}} \caption{#3}
% \end{figure}}
%
% \newcommand{\scalefigcaplabel}[4]{
% \begin{figure}[h]
%  \scalebox{#1}{\includegraphics{#2}}\caption{#3\label{#4}}
% \end{figure}}
%
%-----------------------------------------------------
%Programs
%-----------------------------------------------------
%\newenvironment{prog}[1]{
%\begin{minipage}{5.8 in}
%\begin{center}
%{\sc #1}
%\end{center}
%\begin{itemize}}
%{
%\end{itemize}
%\end{minipage}}

\newcommand{\program}[2]{\fbox{\vspace{2mm}\begin{prog}{#1} #2 \end{prog}\vspace{2mm}}}
%-----------------------------------------------------------

%--------------------------------------------------------------
%--------------------------------------------------------------
%Other - Math
%--------------------------------------------------------------
%--------------------------------------------------------------
\renewcommand{\phi}{\varphi}
\newcommand{\eps}{\epsilon}
\newcommand{\Sum}[1]{\underset{#1}{\Sigma}}
\newcommand{\half}{\ensuremath{\frac{1}{2}}}




%First parameter: index from which counting starts is optional, default is 1; second parameter is the %letter before the property indices
%Example: \begin{properties}{C} or \begin{properties}[5]{C} will start with C6
\newenvironment{properties}[2][0]
{\renewcommand{\theenumi}{#2\arabic{enumi}}
\begin{enumerate} \setcounter{enumi}{#1}}{\end{enumerate}\renewcommand{\theenumi}{\arabic{enumi}}}


\setlength{\parskip}{2mm} \setlength{\parindent}{0mm}
%------------------------------------------------------------
%Making subfigures
%------------------------------------------------------------

%\begin{figure}[h]
%\centering
%\subfigure[caption subfigure 1]{\scalebox{0.2}{\includegraphics{file1..pdf}}\label{fig: label1}}
%\hspace{1cm}
%\subfigure[caption subfigure 2]{
%\scalebox{0.2}{\includegraphics{file2.pdf}}\label{fig: label2}}
%\hspace{1cm}
%\subfigure[caption subfigure 3]{\scalebox{0.2}{\includegraphics{file3.pdf}}\label{fig: label3}}
%\caption{caption for the whole figure\label{fig: whole figure}}
%\end{figure}


%------------------------------------------------------------
%Regular figures
%------------------------------------------------------------

%\begin{figure}[h]
% \scalebox{0.3}{\includegraphics{file.pdf}}\caption{CAPTION \label{fig: label}}
%\end{figure}
